\section{A Heuristic Inertia Score, Not A Proven Law}

\begin{definition}[Heuristic inertia score]
If an analyst chooses a common normalization for the components of
$\Pi(s)$, a \textbf{heuristic inertia score} may be introduced:
\[
  \widehat{\FI}(s) = \alpha \Reward(s) + \beta \Switch(s) + \gamma \Env(s),
\]
with $\alpha,\beta,\gamma > 0$ chosen as modelling weights.
\end{definition}

\begin{remark}
The hat on $\widehat{\FI}$ matters. This expression is a bookkeeping device,
not an established law of behavioural dynamics. It is useful when one wants a
single comparison score, but the manuscript should not pretend that the weights
have been empirically identified or that the linear form is uniquely correct.
\end{remark}

\begin{discussion}
What is defensible here is the \textbf{structure of the scaffold}, not the
claim that the linear combination has already been validated as a universal
formula. Textbook and review sources justify using low-dimensional state
descriptions and stability language. They do not justify claiming that the
human cost of action obeys a known three-term linear law. The score is kept
because it is operationally helpful, but it is explicitly downgraded from
``formal law'' to ``heuristic model.''
\end{discussion}

\begin{remark}
The chapter could have avoided a scalar score entirely and remained purely
ordinal. The reason for retaining $\widehat{\FI}$ is pedagogical economy: later
chapters need a compact way to say that one state is more resistant than
another. The chapter therefore licenses the score as a \emph{comparison aid},
not as a deep empirical result.
\end{remark}
